\subsection{Null and Constant Sequences}
% ---------------------------------------------------------
% Definition --- Constant Sequence
% ---------------------------------------------------------

\begin{definition}[Constant Sequence]
\label{def:constant-sequence}

\textbf{English.}
A sequence $(x_n)$ in $\mathbb{R}$ is constant if all of its terms are equal to a fixed real number.

\textbf{Quantified.}
There exists $c \in \mathbb{R}$ such that for every $n \in \mathbb{N}$,
\[
x_n = c.
\]

\textbf{Fully Quantified Form (FQF).}
\[
\exists c \in \mathbb{R}\;
\forall n \in \mathbb{N}\;
(x_n = c).
\]

\textbf{Predicate.}
\[
\exists c \in \mathbb{R}\;
\mathrm{ConstantSequence}(x_n, c).
\]

\end{definition}

\begin{remark*}[Interpretation]
A constant sequence does not evolve: the “time index” $n$ has no effect on the value.
It represents a completely static process.
\end{remark*}
% ---------------------------------------------------------
% Definition --- Null Sequence
% ---------------------------------------------------------

\begin{definition}[Null Sequence]
\label{def:null-sequence}

\textbf{English.}
A sequence $(x_n)$ in $\mathbb{R}$ is a null sequence if its terms become arbitrarily close to $0$.

\textbf{Quantified.}
For every $\varepsilon > 0$, there exists $N \in \mathbb{N}$ such that for every $n \geq N$,
\[
|x_n| < \varepsilon.
\]

\textbf{Fully Quantified Form (FQF).}
\[
\forall \varepsilon > 0\;
\exists N \in \mathbb{N}\;
\forall n \in \mathbb{N}\;
(n \geq N \rightarrow |x_n| < \varepsilon).
\]

\textbf{Predicate.}
\[
\mathrm{NullSequence}(x_n).
\]

\end{definition}

\begin{remark*}[Interpretation]
A null sequence captures the idea of vanishing: no matter how small a tolerance is chosen,
the sequence eventually remains within that tolerance of zero.
\end{remark*}

% ---------------------------------------------------------
% Theorem --- Constant Sequence Convergence
% ---------------------------------------------------------

\begin{tcolorbox}[
  colback=thmbox,
  colframe=thmborder,
  title={\small\textbf{Theorem (Constant Sequence Convergence)}}
]

\label{thm:constant-sequence-convergence}

\textbf{English.}
Every constant sequence converges to its constant value.

\textbf{Quantified.}
If there exists $c \in \mathbb{R}$ such that for every $n \in \mathbb{N}$,
\[
x_n = c,
\]
then
\[
\lim_{n \to \infty} x_n = c.
\]

\textbf{Fully Quantified Form (FQF).}
\[
\begin{aligned}
&\exists c \in \mathbb{R}\;
\forall n \in \mathbb{N}\;
(x_n = c)
\\
&\qquad\rightarrow
\\
&\forall \varepsilon > 0\;
\exists N \in \mathbb{N}\;
\forall n \in \mathbb{N}\;
(n \geq N \rightarrow |x_n - c| < \varepsilon).
\end{aligned}
\]

\textbf{Predicate.}
\[
\mathrm{ConstantSequence}(x_n, c)
\;\rightarrow\;
\mathrm{ConvergentSequence}(x_n, c).
\]

\end{tcolorbox}

\begin{remark*}[Interpretation]
There is no approximation occurring here: the sequence is already equal to its limit at every index.
\end{remark*}

% ---------------------------------------------------------
% Theorem --- Zero Sequence is Null
% ---------------------------------------------------------

\begin{tcolorbox}[
  colback=thmbox,
  colframe=thmborder,
  title={\small\textbf{Theorem (Zero Sequence is Null)}}
]

\label{thm:zero-sequence-is-null}

\textbf{English.}
The constant zero sequence is a null sequence.

\textbf{Quantified.}
If for every $n \in \mathbb{N}$,
\[
x_n = 0,
\]
then $(x_n)$ is a null sequence.

\textbf{Fully Quantified Form (FQF).}
\[
\begin{aligned}
&\forall n \in \mathbb{N}\;
(x_n = 0)
\\
&\qquad\rightarrow
\\
&\forall \varepsilon > 0\;
\exists N \in \mathbb{N}\;
\forall n \in \mathbb{N}\;
(n \geq N \rightarrow |x_n| < \varepsilon).
\end{aligned}
\]

\textbf{Predicate.}
\[
\mathrm{ConstantSequence}(x_n, 0)
\;\rightarrow\;
\mathrm{NullSequence}(x_n).
\]

\end{tcolorbox}

\begin{remark*}[Interpretation]
The sequence is already exactly at zero, so it trivially satisfies every $\varepsilon$-constraint.
\end{remark*}

% ---------------------------------------------------------
% Theorem --- Constant Null Sequence
% ---------------------------------------------------------

\begin{tcolorbox}[
  colback=thmbox,
  colframe=thmborder,
  title={\small\textbf{Theorem (Constant Null Sequence)}}
]

\label{thm:constant-null-sequence}

\textbf{English.}
A constant sequence is null if and only if its constant value is zero.

\textbf{Quantified.}
If there exists $c \in \mathbb{R}$ such that
\[
x_n = c \quad \text{for every } n \in \mathbb{N},
\]
and $(x_n)$ is a null sequence, then $c = 0$.

\textbf{Fully Quantified Form (FQF).}
\[
\begin{aligned}
&\exists c \in \mathbb{R}\;
\forall n \in \mathbb{N}\;
(x_n = c)
\\
&\land
\\
&\forall \varepsilon > 0\;
\exists N \in \mathbb{N}\;
\forall n \in \mathbb{N}\;
(n \geq N \rightarrow |x_n| < \varepsilon)
\\
&\qquad\rightarrow
\\
&c = 0.
\end{aligned}
\]

\textbf{Predicate.}
\[
\mathrm{ConstantSequence}(x_n, c)
\land
\mathrm{NullSequence}(x_n)
\;\rightarrow\;
c = 0.
\]

\end{tcolorbox}

\begin{remark*}[Interpretation]
A sequence that never changes can only “approach zero” if it is already zero.
\end{remark*}

% ---------------------------------------------------------
% Definition --- Ultimately Constant Sequence
% ---------------------------------------------------------

\begin{definition}[Ultimately Constant Sequence]
\label{def:ultimately-constant-sequence}

\textbf{English.}
A sequence $(x_n)$ in $\mathbb{R}$ is ultimately constant if, after some finite index, all of its terms are equal to one fixed real number.

\textbf{Quantified.}
There exist $N \in \mathbb{N}$ and $c \in \mathbb{R}$ such that for every $n \in \mathbb{N}$,
\[
n \geq N
\quad\Longrightarrow\quad
x_n = c.
\]

\textbf{Fully Quantified Form (FQF).}
\[
\exists N \in \mathbb{N}\;
\exists c \in \mathbb{R}\;
\forall n \in \mathbb{N}\;
(n \geq N \rightarrow x_n = c).
\]

\textbf{Definition predicate reading.}
\[
\mathrm{UltimatelyConstantSequence}(x_n,c).
\]

\end{definition}

\begin{remark*}[Interpretation]
An ultimately constant sequence may change during its initial finite segment, but its tail becomes fixed forever. Thus only the eventual behavior matters.
\end{remark*}

% ---------------------------------------------------------
% Theorem --- Difference from the Limit is Null
% ---------------------------------------------------------

\begin{tcolorbox}[
  colback=thmbox,
  colframe=thmborder,
  title={\small\textbf{Theorem (Difference from the Limit is Null)}}
]

\label{thm:difference-from-limit-is-null}

\textbf{English.}
A sequence converges to $L$ if and only if subtracting $L$ from each term produces a null sequence.

\textbf{Quantified.}
\[
\lim_{n \to \infty} x_n = L
\quad\Longleftrightarrow\quad
(x_n - L) \text{ is a null sequence}.
\]

\textbf{Fully Quantified Form (FQF).}
\[
\begin{aligned}
&\left[
\forall \varepsilon > 0\;
\exists N \in \mathbb{N}\;
\forall n \in \mathbb{N}\;
(n \geq N \rightarrow |x_n - L| < \varepsilon)
\right]
\\
&\qquad\Longleftrightarrow
\\
&\left[
\forall \varepsilon > 0\;
\exists N \in \mathbb{N}\;
\forall n \in \mathbb{N}\;
(n \geq N \rightarrow |(x_n - L)| < \varepsilon)
\right].
\end{aligned}
\]

\textbf{Predicate.}
\[
\mathrm{ConvergentSequence}(x_n, L)
\;\Longleftrightarrow\;
\mathrm{NullSequence}(x_n - L).
\]

\end{tcolorbox}

\begin{remark*}[Interpretation]
This theorem reduces all convergence problems to the special case of convergence to zero,
making null sequences the fundamental building block of limit theory.
\end{remark*}


% =========================================================
% Theorems --- Ultimately Constant Sequences
% =========================================================

% ---------------------------------------------------------
% Theorem --- Ultimately Constant Sequence Convergence
% ---------------------------------------------------------

\begin{tcolorbox}[
  colback=thmbox,
  colframe=thmborder,
  title={\small\textbf{Theorem (Ultimately Constant Sequence Convergence)}}
]
\label{thm:ultimately-constant-sequence-convergence}

\textbf{English.}
Every ultimately constant real sequence converges to its eventual constant value.

\textbf{Quantified.}
If there exist $N \in \mathbb{N}$ and $c \in \mathbb{R}$ such that
\[
n \geq N \;\Longrightarrow\; x_n = c,
\]
then
\[
\lim_{n \to \infty} x_n = c.
\]

\textbf{Fully Quantified Form (FQF).}
\[
\begin{aligned}
&\exists N \in \mathbb{N}\;
\exists c \in \mathbb{R}\;
\forall n \in \mathbb{N}\;
(n \geq N \rightarrow x_n = c)
\\
&\qquad\rightarrow
\\
&\forall \varepsilon > 0\;
\exists N' \in \mathbb{N}\;
\forall n \in \mathbb{N}\;
(n \geq N' \rightarrow |x_n - c| < \varepsilon).
\end{aligned}
\]

\textbf{Predicate.}
\[
\mathrm{UltimatelyConstantSequence}(x_n,c)
\;\rightarrow\;
\mathrm{ConvergentSequence}(x_n,c).
\]

\end{tcolorbox}

\begin{remark*}[Interpretation]
Once the sequence stabilizes, no further approximation is needed—the tail is already equal to the limit.
\end{remark*}

% ---------------------------------------------------------
% Theorem --- Constant Implies Ultimately Constant
% ---------------------------------------------------------

\begin{tcolorbox}[
  colback=thmbox,
  colframe=thmborder,
  title={\small\textbf{Theorem (Constant Implies Ultimately Constant)}}
]
\label{thm:constant-implies-ultimately-constant}

\textbf{English.}
Every constant sequence is ultimately constant.

\textbf{Quantified.}
If there exists $c \in \mathbb{R}$ such that
\[
x_n = c
\qquad
\text{for every } n \in \mathbb{N},
\]
then there exists $N \in \mathbb{N}$ such that
\[
n \geq N \;\Longrightarrow\; x_n = c.
\]

\textbf{Fully Quantified Form (FQF).}
\[
\begin{aligned}
&\exists c \in \mathbb{R}\;
\forall n \in \mathbb{N}\;
(x_n = c)
\\
&\qquad\rightarrow
\\
&\exists N \in \mathbb{N}\;
\exists c \in \mathbb{R}\;
\forall n \in \mathbb{N}\;
(n \geq N \rightarrow x_n = c).
\end{aligned}
\]

\textbf{Predicate.}
\[
\mathrm{ConstantSequence}(x_n,c)
\;\rightarrow\;
\mathrm{UltimatelyConstantSequence}(x_n,c).
\]

\end{tcolorbox}

\begin{remark*}[Interpretation]
A constant sequence is already stabilized from the beginning, so its “eventual” behavior is immediate.
\end{remark*}

% ---------------------------------------------------------
% Theorem --- Ultimately Zero Sequence is Null
% ---------------------------------------------------------

\begin{tcolorbox}[
  colback=thmbox,
  colframe=thmborder,
  title={\small\textbf{Theorem (Ultimately Zero Sequence is Null)}}
]
\label{thm:ultimately-zero-sequence-is-null}

\textbf{English.}
If a sequence is eventually equal to zero, then it is a null sequence.

\textbf{Quantified.}
If there exists $N \in \mathbb{N}$ such that
\[
n \geq N \;\Longrightarrow\; x_n = 0,
\]
then $(x_n)$ is a null sequence.

\textbf{Fully Quantified Form (FQF).}
\[
\begin{aligned}
&\exists N \in \mathbb{N}\;
\forall n \in \mathbb{N}\;
(n \geq N \rightarrow x_n = 0)
\\
&\qquad\rightarrow
\\
&\forall \varepsilon > 0\;
\exists N' \in \mathbb{N}\;
\forall n \in \mathbb{N}\;
(n \geq N' \rightarrow |x_n| < \varepsilon).
\end{aligned}
\]

\textbf{Predicate.}
\[
\mathrm{UltimatelyConstantSequence}(x_n,0)
\;\rightarrow\;
\mathrm{NullSequence}(x_n).
\]

\end{tcolorbox}

\begin{remark*}[Interpretation]
The sequence eventually becomes identically zero, so it trivially satisfies every $\varepsilon$-requirement.
\end{remark*}

% ---------------------------------------------------------
% Theorem --- Ultimately Constant Null Sequence
% ---------------------------------------------------------

\begin{tcolorbox}[
  colback=thmbox,
  colframe=thmborder,
  title={\small\textbf{Theorem (Ultimately Constant Null Sequence)}}
]
\label{thm:ultimately-constant-null-sequence}

\textbf{English.}
An ultimately constant sequence is null if and only if its eventual constant value is zero.

\textbf{Quantified.}
If there exist $N \in \mathbb{N}$ and $c \in \mathbb{R}$ such that
\[
n \geq N \;\Longrightarrow\; x_n = c,
\]
and $(x_n)$ is a null sequence, then
\[
c = 0.
\]

\textbf{Fully Quantified Form (FQF).}
\[
\begin{aligned}
&\exists N \in \mathbb{N}\;
\exists c \in \mathbb{R}\;
\forall n \in \mathbb{N}\;
(n \geq N \rightarrow x_n = c)
\\
&\land
\\
&\forall \varepsilon > 0\;
\exists N' \in \mathbb{N}\;
\forall n \in \mathbb{N}\;
(n \geq N' \rightarrow |x_n| < \varepsilon)
\\
&\qquad\rightarrow
\\
&c = 0.
\end{aligned}
\]

\textbf{Predicate.}
\[
\mathrm{UltimatelyConstantSequence}(x_n,c)
\land
\mathrm{NullSequence}(x_n)
\;\rightarrow\;
c = 0.
\]

\end{tcolorbox}

\begin{remark*}[Interpretation]
A sequence that eventually freezes can only converge to zero if that frozen value is already zero.
\end{remark*}

% ---------------------------------------------------------
% Theorem --- Tail Equality Preserves Convergence
% ---------------------------------------------------------

\begin{tcolorbox}[
  colback=thmbox,
  colframe=thmborder,
  title={\small\textbf{Theorem (Tail Equality Preserves Convergence)}}
]
\label{thm:tail-equality-preserves-convergence}

\textbf{English.}
If two sequences agree from some index onward, then they either both converge or both diverge, and when they converge they have the same limit.

\textbf{Quantified.}
If there exists $N \in \mathbb{N}$ such that
\[
n \geq N \;\Longrightarrow\; x_n = y_n,
\]
then $(x_n)$ converges if and only if $(y_n)$ converges, and if they converge then
\[
\lim_{n \to \infty} x_n = \lim_{n \to \infty} y_n.
\]

\textbf{Fully Quantified Form (FQF).}
\[
\begin{aligned}
&\exists N \in \mathbb{N}\;
\forall n \in \mathbb{N}\;
(n \geq N \rightarrow x_n = y_n)
\\
&\qquad\rightarrow
\\
&\Big(
\big[
\forall \varepsilon > 0\;
\exists N_1 \in \mathbb{N}\;
\forall n \in \mathbb{N}\;
(n \geq N_1 \rightarrow |x_n - L| < \varepsilon)
\big]
\\
&\qquad\qquad\Longleftrightarrow
\big[
\forall \varepsilon > 0\;
\exists N_2 \in \mathbb{N}\;
\forall n \in \mathbb{N}\;
(n \geq N_2 \rightarrow |y_n - L| < \varepsilon)
\big]
\Big).
\end{aligned}
\]

\textbf{Predicate.}
\[
\mathrm{EventuallyEqual}(x_n,y_n)
\;\rightarrow\;
\big(
\mathrm{ConvergentSequence}(x_n,L)
\;\Longleftrightarrow\;
\mathrm{ConvergentSequence}(y_n,L)
\big).
\]

\end{tcolorbox}

\begin{remark*}[Interpretation]
Limits depend only on the tail of a sequence. Finite initial deviations have no influence on convergence behavior.
\end{remark*}











% =========================================================
% Definitions --- Bounded Real Sequences (Full House Style)
% =========================================================

% ---------------------------------------------------------
% Definition --- Sequence Bounded Above
% ---------------------------------------------------------

\begin{definition}[Sequence Bounded Above]
\label{def:sequence-bounded-above}

\textbf{English.}
A real sequence $(x_n)$ is bounded above if all of its terms are less than or equal to one fixed real number.

\textbf{Quantified.}
There exists $M \in \mathbb{R}$ such that for every $n \in \mathbb{N}$,
\[
x_n \leq M.
\]

\textbf{Fully Quantified Form (FQF).}
\[
\exists M \in \mathbb{R}\;
\forall n \in \mathbb{N}\;
(x_n \leq M).
\]

\textbf{Definition predicate reading.}
\[
\mathrm{BoundedAboveSequence}(x_n).
\]

\end{definition}

\begin{remark*}[Negation]
\[
\forall M \in \mathbb{R}\;
\exists n \in \mathbb{N}\;
(M < x_n).
\]
\end{remark*}

\begin{remark*}[Negation predicate reading]
\[
\neg \mathrm{BoundedAboveSequence}(x_n).
\]
\end{remark*}

\begin{remark*}[Failure modes]
The sequence eventually exceeds every proposed upper bound.
\end{remark*}

\begin{remark*}[Failure mode decomposition]
\[
\forall M \in \mathbb{R}\;
\exists n \in \mathbb{N}\;
(x_n > M).
\]
\end{remark*}

\begin{remark*}[Interpretation]
A bounded-above sequence has a ceiling. Failure means it can climb arbitrarily high.
\end{remark*}

% ---------------------------------------------------------
% Definition --- Sequence Bounded Below
% ---------------------------------------------------------

\begin{definition}[Sequence Bounded Below]
\label{def:sequence-bounded-below}

\textbf{English.}
A real sequence $(x_n)$ is bounded below if all of its terms are greater than or equal to one fixed real number.

\textbf{Quantified.}
There exists $m \in \mathbb{R}$ such that for every $n \in \mathbb{N}$,
\[
m \leq x_n.
\]

\textbf{Fully Quantified Form (FQF).}
\[
\exists m \in \mathbb{R}\;
\forall n \in \mathbb{N}\;
(m \leq x_n).
\]

\textbf{Definition predicate reading.}
\[
\mathrm{BoundedBelowSequence}(x_n).
\]

\end{definition}

\begin{remark*}[Negation]
\[
\forall m \in \mathbb{R}\;
\exists n \in \mathbb{N}\;
(x_n < m).
\]
\end{remark*}

\begin{remark*}[Negation predicate reading]
\[
\neg \mathrm{BoundedBelowSequence}(x_n).
\]
\end{remark*}

\begin{remark*}[Failure modes]
The sequence eventually falls below every proposed lower bound.
\end{remark*}

\begin{remark*}[Failure mode decomposition]
\[
\forall m \in \mathbb{R}\;
\exists n \in \mathbb{N}\;
(x_n < m).
\]
\end{remark*}

\begin{remark*}[Interpretation]
A bounded-below sequence has a floor. Failure means it can drop arbitrarily low.
\end{remark*}

% ---------------------------------------------------------
% Definition --- Bounded Sequence
% ---------------------------------------------------------

\begin{definition}[Bounded Sequence]
\label{def:bounded-sequence}

\textbf{English.}
A real sequence $(x_n)$ is bounded if all of its terms remain within one fixed finite distance from zero.

\textbf{Quantified.}
There exists $M \in \mathbb{R}$ with $M > 0$ such that for every $n \in \mathbb{N}$,
\[
|x_n| \leq M.
\]

\textbf{Fully Quantified Form (FQF).}
\[
\exists M \in \mathbb{R}\;
\bigl(
M > 0
\land
\forall n \in \mathbb{N}\;
(|x_n| \leq M)
\bigr).
\]

\textbf{Definition predicate reading.}
\[
\mathrm{BoundedSequence}(x_n).
\]

\end{definition}

\begin{remark*}[Negation]
\[
\forall M \in \mathbb{R}\;
\bigl(
M > 0
\rightarrow
\exists n \in \mathbb{N}\;
(|x_n| > M)
\bigr).
\]
\end{remark*}

\begin{remark*}[Negation predicate reading]
\[
\neg \mathrm{BoundedSequence}(x_n).
\]
\end{remark*}

\begin{remark*}[Failure modes]
The sequence escapes every symmetric bound; its magnitude becomes arbitrarily large.
\end{remark*}

\begin{remark*}[Failure mode decomposition]
\[
\forall M > 0\;
\exists n \in \mathbb{N}\;
(|x_n| > M).
\]
\end{remark*}

\begin{remark*}[Interpretation]
A bounded sequence is trapped in a finite band. Failure means the sequence can grow without bound in magnitude.
\end{remark*}


% =========================================================
% Theorems --- Eventually-Based Tail Formulations for Bounded Sequences
% =========================================================

% ---------------------------------------------------------
% Theorem --- Eventually Bounded Above Tail Formulation
% ---------------------------------------------------------

\begin{tcolorbox}[
  colback=thmbox,
  colframe=thmborder,
  title={\small\textbf{Theorem (Eventually Bounded Above Tail Formulation)}}
]
\label{thm:eventually-bounded-above-tail-formulation}

\textbf{English.}
A real sequence is bounded above if and only if its tail is bounded above.

\textbf{Quantified.}
A sequence $(x_n)$ is bounded above if and only if there exist $N \in \mathbb{N}$ and $M \in \mathbb{R}$ such that
\[
n \geq N
\quad\Longrightarrow\quad
x_n \leq M.
\]

\textbf{Fully Quantified Form (FQF).}
\[
\begin{aligned}
&\left[
\exists M \in \mathbb{R}\;
\forall n \in \mathbb{N}\;
(x_n \leq M)
\right]
\\
&\qquad\Longleftrightarrow
\\
&\left[
\exists N \in \mathbb{N}\;
\exists M \in \mathbb{R}\;
\forall n \in \mathbb{N}\;
(n \geq N \rightarrow x_n \leq M)
\right].
\end{aligned}
\]

\textbf{Predicate reading.}
\[
\mathrm{BoundedAboveSequence}(x_n)
\;\Longleftrightarrow\;
\mathrm{Eventually}(\mathrm{BoundedAboveSequence}(x_n)).
\]

\end{tcolorbox}

\begin{remark*}[Interpretation]
Finite initial terms do not matter for bounded-above behavior. If the tail has a ceiling, the whole sequence has a ceiling after also accounting for the finitely many earlier terms.
\end{remark*}

% ---------------------------------------------------------
% Theorem --- Eventually Bounded Below Tail Formulation
% ---------------------------------------------------------

\begin{tcolorbox}[
  colback=thmbox,
  colframe=thmborder,
  title={\small\textbf{Theorem (Eventually Bounded Below Tail Formulation)}}
]
\label{thm:eventually-bounded-below-tail-formulation}

\textbf{English.}
A real sequence is bounded below if and only if its tail is bounded below.

\textbf{Quantified.}
A sequence $(x_n)$ is bounded below if and only if there exist $N \in \mathbb{N}$ and $m \in \mathbb{R}$ such that
\[
n \geq N
\quad\Longrightarrow\quad
m \leq x_n.
\]

\textbf{Fully Quantified Form (FQF).}
\[
\begin{aligned}
&\left[
\exists m \in \mathbb{R}\;
\forall n \in \mathbb{N}\;
(m \leq x_n)
\right]
\\
&\qquad\Longleftrightarrow
\\
&\left[
\exists N \in \mathbb{N}\;
\exists m \in \mathbb{R}\;
\forall n \in \mathbb{N}\;
(n \geq N \rightarrow m \leq x_n)
\right].
\end{aligned}
\]

\textbf{Predicate reading.}
\[
\mathrm{BoundedBelowSequence}(x_n)
\;\Longleftrightarrow\;
\mathrm{Eventually}(\mathrm{BoundedBelowSequence}(x_n)).
\]

\end{tcolorbox}

\begin{remark*}[Interpretation]
Finite initial terms do not matter for bounded-below behavior. If the tail has a floor, the whole sequence has a floor after also accounting for the finitely many earlier terms.
\end{remark*}

% ---------------------------------------------------------
% Theorem --- Eventually Bounded Tail Formulation
% ---------------------------------------------------------

\begin{tcolorbox}[
  colback=thmbox,
  colframe=thmborder,
  title={\small\textbf{Theorem (Eventually Bounded Tail Formulation)}}
]
\label{thm:eventually-bounded-tail-formulation}

\textbf{English.}
A real sequence is bounded if and only if its tail is bounded.

\textbf{Quantified.}
A sequence $(x_n)$ is bounded if and only if there exist $N \in \mathbb{N}$ and $M > 0$ such that
\[
n \geq N
\quad\Longrightarrow\quad
|x_n| \leq M.
\]

\textbf{Fully Quantified Form (FQF).}
\[
\begin{aligned}
&\left[
\exists M \in \mathbb{R}\;
\bigl(
M > 0
\land
\forall n \in \mathbb{N}\;
(|x_n| \leq M)
\bigr)
\right]
\\
&\qquad\Longleftrightarrow
\\
&\left[
\exists N \in \mathbb{N}\;
\exists M \in \mathbb{R}\;
\bigl(
M > 0
\land
\forall n \in \mathbb{N}\;
(n \geq N \rightarrow |x_n| \leq M)
\bigr)
\right].
\end{aligned}
\]

\textbf{Predicate reading.}
\[
\mathrm{BoundedSequence}(x_n)
\;\Longleftrightarrow\;
\mathrm{Eventually}(\mathrm{BoundedSequence}(x_n)).
\]

\end{tcolorbox}

\begin{remark*}[Interpretation]
Boundedness is a tail property up to finitely many exceptions. A finite number of large early terms can be absorbed into a larger bound.
\end{remark*}

% =========================================================
% Theorems --- Relationships Between Bounded Sequence Definitions
% =========================================================

% ---------------------------------------------------------
% Theorem --- Bounded Sequence iff Bounded Above and Below
% ---------------------------------------------------------

\begin{tcolorbox}[
  colback=thmbox,
  colframe=thmborder,
  title={\small\textbf{Theorem (Bounded Sequence iff Bounded Above and Below)}}
]
\label{thm:bounded-sequence-iff-bounded-above-and-below}

\textbf{English.}
A real sequence is bounded if and only if it is both bounded above and bounded below.

\textbf{Quantified.}
A real sequence $(x_n)$ is bounded if and only if there exist real numbers $m,M \in \mathbb{R}$ such that
\[
m \leq x_n \leq M
\qquad
\text{for every } n \in \mathbb{N}.
\]

\textbf{Fully Quantified Form (FQF).}
\[
\begin{aligned}
&\left[
\exists K \in \mathbb{R}\;
\bigl(
K > 0
\land
\forall n \in \mathbb{N}\;
(|x_n| \leq K)
\bigr)
\right]
\\
&\qquad\Longleftrightarrow
\\
&\left[
\exists m \in \mathbb{R}\;
\exists M \in \mathbb{R}\;
\forall n \in \mathbb{N}\;
(m \leq x_n \land x_n \leq M)
\right].
\end{aligned}
\]

\textbf{Predicate reading.}
\[
\mathrm{BoundedSequence}(x_n)
\;\Longleftrightarrow\;
\bigl(
\mathrm{BoundedBelowSequence}(x_n)
\land
\mathrm{BoundedAboveSequence}(x_n)
\bigr).
\]

\textbf{Contrapositive predicate reading.}
\[
\neg\bigl(
\mathrm{BoundedBelowSequence}(x_n)
\land
\mathrm{BoundedAboveSequence}(x_n)
\bigr)
\;\Longleftrightarrow\;
\neg\mathrm{BoundedSequence}(x_n).
\]

\end{tcolorbox}

\begin{remark*}[Interpretation]
Boundedness means the sequence is trapped between two finite horizontal barriers.
The symmetric absolute-value bound $|x_n|\leq K$ is equivalent to having both a floor and a ceiling.
\end{remark*}

% ---------------------------------------------------------
% Theorem --- Absolute Bound Gives Upper and Lower Bounds
% ---------------------------------------------------------

\begin{tcolorbox}[
  colback=thmbox,
  colframe=thmborder,
  title={\small\textbf{Theorem (Absolute Bound Gives Upper and Lower Bounds)}}
]
\label{thm:absolute-bound-gives-upper-and-lower-bounds}

\textbf{English.}
If every term of a real sequence lies within distance $K$ of zero, then $K$ is an upper bound and $-K$ is a lower bound.

\textbf{Quantified.}
If $K>0$ and
\[
|x_n|\leq K
\qquad
\text{for every } n\in\mathbb{N},
\]
then
\[
-K\leq x_n\leq K
\qquad
\text{for every } n\in\mathbb{N}.
\]

\textbf{Fully Quantified Form (FQF).}
\[
\begin{aligned}
&\forall K \in \mathbb{R}\;
\Bigl(
K > 0
\land
\forall n \in \mathbb{N}\;
(|x_n| \leq K)
\Bigr)
\\
&\qquad\rightarrow
\\
&\forall n \in \mathbb{N}\;
(-K \leq x_n \land x_n \leq K).
\end{aligned}
\]

\textbf{Predicate reading.}
\[
\mathrm{BoundedSequence}(x_n)
\;\rightarrow\;
\bigl(
\mathrm{BoundedBelowSequence}(x_n)
\land
\mathrm{BoundedAboveSequence}(x_n)
\bigr).
\]

\textbf{Contrapositive predicate reading.}
\[
\neg\bigl(
\mathrm{BoundedBelowSequence}(x_n)
\land
\mathrm{BoundedAboveSequence}(x_n)
\bigr)
\;\rightarrow\;
\neg\mathrm{BoundedSequence}(x_n).
\]

\end{tcolorbox}

\begin{remark*}[Interpretation]
An absolute-value bound gives a symmetric trap around zero. The same number that controls the magnitude also supplies explicit upper and lower barriers.
\end{remark*}

% ---------------------------------------------------------
% Theorem --- Upper and Lower Bounds Give an Absolute Bound
% ---------------------------------------------------------

\begin{tcolorbox}[
  colback=thmbox,
  colframe=thmborder,
  title={\small\textbf{Theorem (Upper and Lower Bounds Give an Absolute Bound)}}
]
\label{thm:upper-and-lower-bounds-give-absolute-bound}

\textbf{English.}
If a real sequence has both a lower bound and an upper bound, then it has an absolute-value bound.

\textbf{Quantified.}
If there exist $m,M\in\mathbb{R}$ such that
\[
m\leq x_n\leq M
\qquad
\text{for every } n\in\mathbb{N},
\]
then there exists $K>0$ such that
\[
|x_n|\leq K
\qquad
\text{for every } n\in\mathbb{N}.
\]

\textbf{Fully Quantified Form (FQF).}
\[
\begin{aligned}
&\left[
\exists m \in \mathbb{R}\;
\exists M \in \mathbb{R}\;
\forall n \in \mathbb{N}\;
(m \leq x_n \land x_n \leq M)
\right]
\\
&\qquad\rightarrow
\\
&\left[
\exists K \in \mathbb{R}\;
\bigl(
K > 0
\land
\forall n \in \mathbb{N}\;
(|x_n| \leq K)
\bigr)
\right].
\end{aligned}
\]

\textbf{Predicate reading.}
\[
\bigl(
\mathrm{BoundedBelowSequence}(x_n)
\land
\mathrm{BoundedAboveSequence}(x_n)
\bigr)
\;\rightarrow\;
\mathrm{BoundedSequence}(x_n).
\]

\textbf{Contrapositive predicate reading.}
\[
\neg\mathrm{BoundedSequence}(x_n)
\;\rightarrow\;
\neg\bigl(
\mathrm{BoundedBelowSequence}(x_n)
\land
\mathrm{BoundedAboveSequence}(x_n)
\bigr).
\]

\end{tcolorbox}

\begin{remark*}[Interpretation]
Two one-sided barriers can always be replaced by one symmetric barrier around zero.
A standard choice is any positive $K$ satisfying both $K\geq |m|$ and $K\geq |M|$.
\end{remark*}